\section*{अध्याय \ref{ch:g10:electron-shells} --- इलेक्ट्रॉन कोश और आवर्त सारणी}

\begin{solution}{exo:g10:electron-shells:1}
C: $1\mathrm{s}^2\,2\mathrm{s}^2\,2\mathrm{p}^2$। N:
$1\mathrm{s}^2\,2\mathrm{s}^2\,2\mathrm{p}^3$। Mg:
$1\mathrm{s}^2\,2\mathrm{s}^2\,2\mathrm{p}^6\,3\mathrm{s}^2$।
\end{solution}

\begin{solution}{exo:g10:electron-shells:2}
कार्बन 4, नाइट्रोजन 5, मैग्नीशियम 2।
\end{solution}

\begin{solution}{exo:g10:electron-shells:3}
बाहरी \omterm{def:g10:electron-shells:configuration}{कोश} $n = 3$: \omterm{def:g9:periodic-table-first-look:period}{आवर्त} 3। $2 + 3 = 5$ \omterm{def:g10:electron-shells:valence}{संयोजकता इलेक्ट्रॉन}: \omterm{def:g9:periodic-table-first-look:period}{समूह} 15
(फ़ॉस्फ़ोरस)।
\end{solution}

\begin{solution}{exo:g10:electron-shells:4}
\omterm{def:g10:electron-shells:family}{क्षार धातुएँ} (लिथियम, सोडियम, पोटैशियम); \omterm{def:g10:electron-shells:family}{हैलोजन} (फ़्लुओरीन,
क्लोरीन, ब्रोमीन, आयोडीन); \omterm{def:g10:electron-shells:family}{उत्कृष्ट गैसें} (हीलियम, नियॉन, आर्गन)।
\end{solution}

\begin{solution}{exo:g10:electron-shells:5}
\omterm{def:g7:atoms-and-molecules:atom}{परमाणु} इतने \omterm{def:g9:inside-the-atom:nucleus}{इलेक्ट्रॉन} लेने या खोने की प्रवृत्ति रखते हैं कि उनके बाहरी \omterm{def:g10:electron-shells:configuration}{कोश} में
निकटतम \omterm{def:g10:electron-shells:family}{उत्कृष्ट गैस} की तरह आठ \omterm{def:g9:inside-the-atom:nucleus}{इलेक्ट्रॉन} हो जाएँ।
\end{solution}

\begin{solution}{exo:g10:electron-shells:6}
K, $2.8.8.1$, एक खोता है: \ce{K+},
$1\mathrm{s}^2\,2\mathrm{s}^2\,2\mathrm{p}^6\,3\mathrm{s}^2\,3\mathrm{p}^6$
(आर्गन)। F, $2.7$, एक लेता है: \ce{F-},
$1\mathrm{s}^2\,2\mathrm{s}^2\,2\mathrm{p}^6$ (नियॉन)।
\end{solution}

\begin{solution}{exo:g10:electron-shells:7}
$1\mathrm{s}^2\,2\mathrm{s}^2\,2\mathrm{p}^6\,3\mathrm{s}^2$: $Z = 12$,
मैग्नीशियम।
\end{solution}

\begin{solution}{exo:g10:electron-shells:8}
\ce{Mg^{2+}} और \ce{Cl-}: \ce{MgCl2}। \ce{Al^{3+}} और \ce{O^{2-}}:
\ce{Al2O3}।
\end{solution}

\begin{solution}{exo:g10:electron-shells:9}
सल्फ़र (6 \omterm{def:g10:electron-shells:valence}{संयोजकता इलेक्ट्रॉन}, $2.8.6$), उसी समूह, 16, में।
\end{solution}

\begin{solution}{exo:g10:electron-shells:10}
उसका बाहरी \omterm{def:g10:electron-shells:configuration}{कोश} पहले से भरा है (आठ \omterm{def:g9:inside-the-atom:nucleus}{इलेक्ट्रॉन}): \omterm{def:g9:inside-the-atom:nucleus}{इलेक्ट्रॉन} खोने या लेने से उसे
कुछ नहीं मिलता।
\end{solution}

\begin{solution}{exo:g10:electron-shells:11}
आर्गन में 18 \omterm{def:g9:inside-the-atom:nucleus}{इलेक्ट्रॉन} हैं; \ce{X^{2+}} ने 2 खोए हैं, इसलिए X में 20 हैं: कैल्सियम।
\end{solution}

\begin{solution}{exo:g10:electron-shells:12}
हीलियम, $1\mathrm{s}^2$, का बाहरी \omterm{def:g10:electron-shells:configuration}{कोश} भरा है (\omterm{def:g10:electron-shells:configuration}{कोश} 1 में केवल दो
\omterm{def:g9:inside-the-atom:nucleus}{इलेक्ट्रॉन} आते हैं), दूसरी \omterm{def:g10:electron-shells:family}{उत्कृष्ट गैसों} की तरह; वह \omterm{def:g9:periodic-table-first-look:period}{समूह} 2 की \omterm{def:g9:periodic-table-first-look:metal}{धातुओं} जैसा
व्यवहार नहीं करती।
\end{solution}

\begin{solution}{exo:g10:electron-shells:13}
उसका \omterm{def:g10:electron-shells:valence}{संयोजकता इलेक्ट्रॉन} \omterm{def:g10:electron-shells:configuration}{कोश} 4 में है, \omterm{def:g10:electron-shells:configuration}{कोश} 3 वाले सोडियम के \omterm{def:g9:inside-the-atom:nucleus}{इलेक्ट्रॉन} से \omterm{def:g9:inside-the-atom:nucleus}{नाभिक}
से अधिक दूर: ढीली पकड़ के कारण वह अधिक आसानी से निकल जाता है।
\end{solution}

\begin{solution}{exo:g10:electron-shells:14}
चारों में 18 \omterm{def:g9:inside-the-atom:nucleus}{इलेक्ट्रॉन} हैं और आर्गन का विन्यास:
\[ 1\mathrm{s}^2\,2\mathrm{s}^2\,2\mathrm{p}^6\,3\mathrm{s}^2\,3\mathrm{p}^6. \]
\end{solution}

\begin{solution}{exo:g10:electron-shells:15}
अष्टक पाने के लिए वह 3 \omterm{def:g9:inside-the-atom:nucleus}{इलेक्ट्रॉन} लेता है: उसमें 5 \omterm{def:g10:electron-shells:valence}{संयोजकता इलेक्ट्रॉन} हैं, \omterm{def:g9:periodic-table-first-look:period}{समूह}
15। \omterm{def:g9:periodic-table-first-look:period}{आवर्त} 3 में यह फ़ॉस्फ़ोरस है:
\[ 1\mathrm{s}^2\,2\mathrm{s}^2\,2\mathrm{p}^6\,3\mathrm{s}^2\,3\mathrm{p}^3. \]
\end{solution}

\begin{solution}{pb:g10:electron-shells:1}
\textbf{1.} ऐलुमिनियम: 13 \omterm{def:g9:inside-the-atom:proton}{प्रोटॉन}, 13 \omterm{def:g9:inside-the-atom:nucleus}{इलेक्ट्रॉन}। ऑक्सीजन: 8 \omterm{def:g9:inside-the-atom:proton}{प्रोटॉन}, 8
\omterm{def:g9:inside-the-atom:nucleus}{इलेक्ट्रॉन}।

\textbf{2.} Al: $1\mathrm{s}^2\,2\mathrm{s}^2\,2\mathrm{p}^6\,3\mathrm{s}^2\,3\mathrm{p}^1$।
O: $1\mathrm{s}^2\,2\mathrm{s}^2\,2\mathrm{p}^4$।

\textbf{3.} Al: 3 \omterm{def:g10:electron-shells:valence}{संयोजकता इलेक्ट्रॉन}, \omterm{def:g9:periodic-table-first-look:period}{आवर्त} 3, \omterm{def:g9:periodic-table-first-look:period}{समूह} 13। O: 6, \omterm{def:g9:periodic-table-first-look:period}{आवर्त} 2,
\omterm{def:g9:periodic-table-first-look:period}{समूह} 16।

\textbf{4.} ऐलुमिनियम \omterm{def:g9:periodic-table-first-look:metal}{धातु} है, ऑक्सीजन \omterm{def:g9:periodic-table-first-look:metal}{अधातु}।

\textbf{5.} \ce{Al^{3+}}: $1\mathrm{s}^2\,2\mathrm{s}^2\,2\mathrm{p}^6$।

\textbf{6.} \ce{O^{2-}}: $1\mathrm{s}^2\,2\mathrm{s}^2\,2\mathrm{p}^6$।

\textbf{7.} नियॉन।

\textbf{8.} ऐलुमिनियम 3 खोता है, ऑक्सीजन 2 लेती है।

\textbf{9.} $2 \times (+3) + 3 \times (-2) = 0$: \ce{Al2O3}।

\textbf{10.} खोए: $2 \times 3 = 6$। लिए: $3 \times 2 = 6$।

\textbf{11.} \omterm{def:g9:inside-the-atom:nucleus}{इलेक्ट्रॉन} न बनते हैं, न नष्ट होते हैं: ऐलुमिनियम द्वारा खोए गए ठीक
वही हैं जो ऑक्सीजन ने लिए, और यौगिक
उदासीन है।

\textbf{12.} उस पर वही आवेश है, $+3$: क्रिस्टल के आवेश अब भी
संतुलित रहते हैं।

\textbf{13.} $2 \times 27 + 3 \times 16 = 102$ \omterm{def:g9:inside-the-atom:proton}{न्यूक्लिऑन}।

\textbf{14.} $102 \times 1.67 \times 10^{-27} = \qty{1.70e-25}{kg}$।

\textbf{15.} $10^{-3} / 1.70 \times 10^{-25} \approx 5.87 \times 10^{21}$
सूत्र-इकाइयाँ।

\textbf{16.} $2 \times 5.87 \times 10^{21} = 1.17 \times 10^{22}$
\ce{Al^{3+}} \omterm{def:g9:ions:ion}{आयन}; $3 \times 5.87 \times 10^{21} = 1.76 \times 10^{22}$
\ce{O^{2-}} \omterm{def:g9:ions:ion}{आयन}।

\textbf{17.} \omterm{def:g9:inside-the-atom:nucleus}{इलेक्ट्रॉन} \omterm{def:g9:inside-the-atom:proton}{न्यूक्लिऑन} से लगभग 1836 गुना हल्का है; एक सूत्र-इकाई के
50 \omterm{def:g9:inside-the-atom:nucleus}{इलेक्ट्रॉनों} का भार एक \omterm{def:g9:inside-the-atom:proton}{न्यूक्लिऑन} के भार के तीन प्रतिशत से भी
कम है।

\textbf{18.} \qty{1}{g} कुरुंड के लिए $6 \times 5.87 \times 10^{21} \approx 3.5 \times 10^{22}$
\omterm{def:g9:inside-the-atom:nucleus}{इलेक्ट्रॉन} स्थानांतरित हुए।
\end{solution}
